\documentclass{article}
\usepackage{../math_template}

\author{Jonathan Kasongo}
\title{Mock Edexcel Further Pure 2}

\begin{document}
\maketitle

I wrote this mock for myself to prepare, the questions are from the
FP2 book. It is meant to be harder than the usual exam difficulty.
Unlike the real exam the questions are not ordered by difficulty.
The time limit is 1hr 10m. Good luck.

\begin{problem}{Problem 1}
For non-negative integer $n$, and real $a>0$, let
$$I_n = \int_0^a x^n (a-x)^\frac13 \text{d}x.$$
\vspace{0.2cm}

(a) Prove that $I_n = \frac{3an}{3n+4}I_{n-1}$, for $n\geq 1$ [6].
\vspace{0.2cm}

(b) Given that $I_2 = \frac{27}{49} a^\frac43$, find the value of $a$ [6].
\end{problem}

\begin{problem}{Problem 2}
The matrix
$$
\mathbf{A} = \begin{bmatrix}
3 & 2 & 4 \\
2 & 0 & 2 \\
4 & 2 & k
\end{bmatrix}
$$
has eigenvector $\langle 0, 2, -1 \rangle$.
\vspace{0.2cm}

(a) (i) Find the eigenvalue of $\mathbf{A}$ corresponding to this given
eigenvector. \vspace{0.2cm}

(ii) Find $k$. [4] \vspace{0.2cm}

Another eigenvalue of $\mathbf{A}$ is 8.
It is known that $\mathbf{A}$ has a repeated eigenvalue.
\vspace{0.2cm}

(b) (i) Prove which of the eigenvalues is repeated.
\vspace{0.2cm}

(ii) Using each eigenvalue, find two more eigenvectors of $\mathbf{A}$ such
that both are linearly independent to the eigenvector given in the question.
[8]
\end{problem}

\begin{problem}{Problem 3}
At a family party, the caterers made $n$ similar mini cupcakes.
$n$ was chosen so that the cupcakes could be distributed equally between
the 18 people expected to attend. \vspace{0.2cm}

Due to an illness only 14 people attended the party. To avoid arguments,
the cupcakes were shared evenly amongst these 14 people.
2 cupcakes were left over. \vspace{0.2cm}

(a) Formulate a congruence equation to represent this information.
[2] \vspace{0.2cm}

(b) Given that $n<200$, find two possible values of $n$.
[5]
\end{problem}

\begin{problem}{Problem 4}
$S_4$ is the group of all possible permutations of 4 elements under the
operation of composition. \vspace{0.2cm}

(a) Show that the permutation
$\left(\begin{array}{llll}1 & 2 & 3 & 4 \\ 1 & 4 & 2 & 3\end{array}\right)$
generates a subgroup of $S_4$ of order 3.
[3]
\vspace{0.2cm}

(b) Find a subgroup of $S_4$ of order 2.
[2]
\end{problem}

\begin{problem}{Problem 5}
The transformation $T$ from the $z$-plane, where $z=x+\mathrm{i} y$, to the
$w$-plane where $w=u+\mathrm{i} v$, is given by $w=z+\frac{4}{z}, z \neq 0$.
\vspace{0.2cm}

Show that the transformation $T$ maps the points on a circle $|z|=2$ to
points in the interval $[-k, k]$ on the real axis. State the value of the
constant $k$.

[7]
\end{problem}

\begin{problem}{Problem 6}
The curve with polar equation $r=\mathrm{e}^\theta$, between where
$\theta=0$ and $\theta=\frac{\pi}{2}$, is rotated completely about the
vertical line $\theta= \pm \frac{\pi}{2}$, forming a surface of revolution.
\vspace{0.2cm}

Find the exact area of this surface.
[6]
\end{problem}

\begin{problem}{Problem 7}
The rigid symmetries of a regular hexagon $A B C D E F$ form a group under
the operation of composition. This group contains the element $p$
representing a reflection in the line through $A$ and $D$, and the element
$q$ representing a rotation through $60^{\circ}$ anticlockwise about the
centre of the hexagon. \vspace{0.2cm}

(a) Write down the order of $p$ and the order of $q$. [2] \vspace{0.2cm}

(b) Construct a Caley table for the group generated by $p$. [4]
\vspace{0.2cm}

(c) Describe the effect of the transformation $q^2$, and write down the
elements of the subgroup generated by $q^2$ in terms of $q$. [3]
\end{problem}

\end{document}
